%% ============================================================
%% preamble/commands.tex — semantic inline macros + apparatus + algorithms
%%
%% Purpose      : the inline vocabulary sections may use, the keyword/JEL
%%                apparatus macros (\Keywords, \JEL, \MSC), and the
%%                algorithm environment loader (module: algorithms).
%% Dependencies : etoolbox; code.tex owns \code when the code module is
%%                on (this file supplies a \texttt fallback otherwise).
%% Provenance   : \keyterm from the SSRN cluster; \Keywords/\JEL promoted
%%                from datacenter/moratorium commands.tex to read the
%%                generated apparatus macros.
%% ============================================================

\usepackage{xspace}   % for number-macros (\newcommand{\nX}{12{,}345\xspace})

% --- Number macros (single source of truth) ---------------------------
% The house convention (docs/guides/STYLE-PAPER.md): define every key
% statistic ONCE here and cite the macro, so prose and captions never
% drift. \xspace eats the right trailing space. The sample uses these to
% demonstrate the pattern; replace them with your paper's real figures:
\newcommand{\nEngines}{three\xspace}      % pdflatex, xelatex, lualatex
\newcommand{\nProfiles}{four\xspace}      % libertinus, newtx, lmodern, plex
\newcommand{\nVenues}{three\xspace}       % preprint, arxiv, ssrn
% e.g. real statistics would look like:
%   \newcommand{\nSamples}{12{,}345\xspace}
%   \newcommand{\accuracy}{92.4\%\xspace}

% --- Inline semantics -------------------------------------------------
\newcommand{\term}[1]{\emph{#1}}               % first use of a technical term
\newcommand{\keyterm}[1]{\textbf{\color{primary}#1}}  % defined term at its site
\newcommand{\work}[1]{\emph{#1}}               % titles of works
\newcommand{\foreignphrase}[1]{\emph{#1}}      % non-English phrase

% Inline code — a robust \texttt wrapper (works in tables, captions, and
% footnotes, unlike \lstinline). Escape specials the standard LaTeX way
% inside it: \_ \# \% \{ \} and \textbackslash for a literal backslash.
% Block code with highlighting is the codelisting environment (code.tex).
\newcommand{\code}[1]{\texttt{#1}}

% --- Apparatus (keywords / JEL / MSC) --------------------------------
% Rendered on the abstract page by frontmatter/titleblock; each reads the
% generated macro and prints only when the matching flag is set.
\newcommand{\KeywordsLine}{%
  \ifPaperHasKeywords
    \noindent{\small\textbf{\color{pt-apparatus-label}Keywords:}\ \PaperKeywords\par}%
  \fi}
\newcommand{\JELLine}{%
  \ifPaperHasJEL
    \noindent{\small\textbf{\color{pt-apparatus-label}JEL classification:}\ \PaperJEL\par}%
  \fi}
\newcommand{\MSCLine}{%
  \ifPaperHasMSC
    \noindent{\small\textbf{\color{pt-apparatus-label}MSC:}\ \PaperMSC\par}%
  \fi}

% --- Algorithms (module: algorithms) ---------------------------------
\ifPaperModuleAlgo
  \usepackage{algorithm}
  \usepackage{algpseudocodex}   % modern algorithmicx; falls back cleanly
\fi

% --- Block-quote attribution ------------------------------------------
\newcommand{\attribution}[1]{%
  \par\nopagebreak\smallskip
  {\raggedleft\normalfont\color{text-muted}\textemdash\,#1\par}}

% Block quotes sit a touch smaller than body copy (a common convention;
% the size change also signals "quotation" in grayscale where no color
% cue exists).
\AtBeginEnvironment{quotation}{\small}
